\documentclass[a4paper,11pt]{article}
\title{\textbf{}}
\author{}	
\date{}

\usepackage[spanish]{babel} % espanol
\usepackage{graphicx} % graficos
\usepackage{listings} % para el codigo
\usepackage{vmargin} % margenes 

\begin{document}

\begin{titlepage}
\begin{center}

\vspace*{-1in}
\begin{figure}[htb]
\begin{center}
\includegraphics[width=4cm]{logo.png}
\end{center}
\end{figure}
	
\begin{huge}
\textbf{Universidad de Buenos Aires\\
Facultad de Ingenier\'ia \\}
\vspace*{0.2in}
\end{huge}

\begin{LARGE}
\textbf{66.20 Organizaci\'on de computadoras \\}
\vspace*{1in}
\end{LARGE}

\begin{LARGE}
\textbf{Trabajo pr\'actico 1:
Conjunto de \\instrucciones MIPS \\
%\vspace*{0.2in}
}
\end{LARGE}
\begin{LARGE}
\textbf{1$^{er}$ cuatrimestre de 2015} \\
\end{LARGE}

\vspace*{1in}

\begin{Large}
\begin{tabular}{ l l l }
   Paula Saffioti & Padr\'on: 92001 & Email: paula.saffioti@gmail.com \\
   Kaoru Heanna & Padr\'on: 91891 & Email: kaoru.heanna@gmail.com \\
   Juli\'an Scialabba & Padr\'on: 92181 & Email: julian.scialabba@gmail.com \\
 \end{tabular}\\
\end{Large}

\end{center}
\end{titlepage}

\section{Documentaci\'on}
El objetivo del trabajo pr\'actico es desarrollar una versi\'on del comando tac de UNIX utilizando el conjunto de instrucciones MIPS. 
Ya contab\'amos con el comando implementado en lenguaje C y deb\'iamos traducirlo a assembly MIPS. Para hacer esto optamos por dividir el c\'odigo en funciones peque\~nas, con el objetivo de poder implementarlas y testearlas independientemente.
El algoritmo base del programa que se resume en la siguiente lista de pasos:
\begin{enumerate}
	\item Abrir un archivo pasado por par\'ametro o recibir un stream de datos desde stdin.
	\item Leer el archivo o el stream l\'inea por l\'inea hasta el final y guardarlo en un array. El tama\~no de las l\'ineas debe ser din\'amico.
	\item Recorrer el array en sentido inverso imprimiendo por stdout el contenido de cada l\'inea.
	\item Cerrar el archivo si corresponde.
\end{enumerate}

El paso 1 fue implementado en C. Para el caso de los archivos lo resolvimos con la funci\'on fopen. Como resultado de este paso, nos qued\'o un puntero al archivo para operar con \'el. En caso de alg\'un error, lo lanzamos por stderr.

Los pasos 2 y 3 fueron realizados \'integramente en assembly MIPS. Para el manejo de la memoria din\'amica utilizamos las funciones ya implementadas mymalloc y myfree. Tuvimos que modificar el c\'odigo que utilizaba realloc para que tambi\'en pueda hacer uso de mymalloc.


El paso 4 para cerrar el archivo nuevamente fue implementado en C con la funci\'on fclose. 


El \'unico inconveniente durante el desarrollo de trabajo pr\'actico lo encontramos con la implementaci\'on de fgets. El problema era que haciamos la llamada al syscall encargado de la lectura de archivos (el 14) con los par\'ametros necesarios pero no nos devolv\'ia el bloque le\'ido. Tras mucho intentar, consultamos en el grupo de la materia esta cuesti\'on y al no obtener respuesta lo solucionamos llamando a fgets desde el c\'odigo assembly.

\section{Comandos}
Para obtener el ejecutable creamos un makefile. Por lo que para compilar el programa solo hace falta hacer:

\lstset{language=bash, breaklines=true, basicstyle=\normalsize}
\begin{lstlisting}
$ make
\end{lstlisting}

El makefile es el siguiente:
\begin{lstlisting}
PROG  = tp1
CFLAGS = -g -Wall
DSRC = src/
all: $(PROG)

deleteOld:
	rm $(PROG)

tp1: $(DSRC)main.c $(DSRC)printHelp.S $(DSRC)printVersion.S $(DSRC)isEndOfLine.S $(DSRC)strLength.S $(DSRC)concatBuffer.S $(DSRC)tacFile.S $(DSRC)printLines.S $(DSRC)resizeArrayLines.S  $(DSRC)mergeStrings.S $(DSRC)mymalloc.S $(DSRC)storeNewLine.S $(DSRC)readFile.S
	$(CC) $(CFLAGS) -o $@ $>

clean:
	rm -f $(PROG) *.so *.o *.a *.core
	
permissions:
	chmod +x $(PROG)

\end{lstlisting}


\section{Corridas de prueba}
Para corroborar el funcionamiento del programa cont\'abamos con una serie de archivos de prueba y con el diccionario de palabras ubicado en /usr/share/dict/words. 
\\
\\
\underline{Salidas de ejemplo:}\\
\textbf{return.txt:}\\
\emph{of -1 is returned and errno is set to indicate the error.
Upon successful completion a value of 0 is returned. Otherwise, a value}
\\\\
\textbf{basic.txt}\\
\emph{4. D\\
3. C\\
2. B\\
1. A\\}\\
\textbf{return.txt basic.txt}\\
\emph{of -1 is returned and errno is set to indicate the error.\\
Upon successful completion a value of 0 is returned. Otherwise, a value\\
4. D\\
3. C\\
2. B\\
1. A}\\

Para poder hacer todas las pruebas de manera r\'apida escribimos un script que compilaba el programa, luego lo ejecutaba para los archivos mencionados d\'andolos vuelta dos veces y finalmente verificaba que el archivo resultante fuera igual al original con md5sum. El script es el siguiente:
\lstset{language=bash, breaklines=true, basicstyle=\normalsize}
\begin{lstlisting}
testFile() {
	echo -e "\nTEST "$1".txt"
	./tp1 test-files/$1.txt | ./tp1 > output/$1.txt
	res1=`cat test-files/$1.txt | md5`
	res2=`cat output/$1.txt | md5`
	if [ $res1 == $res2 ]; 
	then
		echo -e "OK\n"
	else
		echo -e "ERROR\n"
	fi
}

echo Compilando...
make

testFile empty
testFile basic
testFile empty-lines
testFile large-file
testFile return
testFile status
\end{lstlisting}
\underline{MD5 de los archivos de prueba procesados dos veces:}\\
\textbf{empty.txt:} d41d8cd98f00b204e9800998ecf8427e\\
\textbf{basic:} f82584604433de83165f0227f4f06c2c\\
\textbf{empty-lines.txt:} 0c060d8e86deedcb0f9cd89f6a3f9c93\\
\textbf{large-file.txt:} 8b2a59a337c384be9b640a38a377eef1\\
\textbf{return.txt:} 6777acdeb99f2fb0ed9e14fc82899e82\\
\textbf{status.txt:} 964ccaae773b2f9044dc549ed3636ab6\\
\textbf{words:} cbbcded3dc3b61ad50c4e36f79c37084

\section{C\'odigo fuente}
\subsection{main.c}

\lstset{language=C, breaklines=true, basicstyle=\normalsize}
\begin{lstlisting}
#include <stdio.h>
#include <stdlib.h>
#include <string.h>

extern void printHelp();
extern void printVersion();
extern int tacFile(FILE* fp);

int main(int argc, char** argv) {
    
    if ((argc == 2) && ((strcmp(argv[1], "-h") == 0) || (strcmp(argv[1], "--help") == 0))) {
        printHelp();
        return (EXIT_SUCCESS);
    }

    if ((argc == 2) && ((strcmp(argv[1], "-V") == 0) || (strcmp(argv[1], "--version") == 0))) {
        printVersion();
        return (EXIT_SUCCESS);
    }

    int result;

    if (argc < 2) { //no tengo archivo de entrada, uso standard input
        result = tacFile(stdin);
        return (result);
    }

    int i;
    for (i = 1; i < argc; i++) {
        FILE *fp;
        fp = fopen(argv[i], "r");
        if (fp == NULL) {
            fprintf(stderr, "%s", argv[i]);
            fprintf(stderr, ": nombre de archivo o comando inválido.\n");
            return (EXIT_FAILURE);
        }
        tacFile(fp);
        fclose(fp);
    }
    return (EXIT_SUCCESS);
}
\end{lstlisting}

\subsection{printHelp.S}

\lstset{language=[x86masm]Assembler, breaklines=true, basicstyle=\normalsize}
\begin{lstlisting}
#include <mips/regdef.h>
#include <sys/syscall.h>

	.text
	.align	2
	.globl	printHelp
	.ent	printHelp
printHelp:
	.frame	$fp,40,ra	#4SRA + 2LTA + 4ABA
	.set	noreorder
	.cpload	t9
	.set	reorder
	subu 	sp,sp,40
	.cprestore 24
	sw	$fp,28(sp) 
	sw	ra,32(sp)
	move 	$fp, sp
	la 	a0, msg
	la	t9, strLength
	jal	ra, t9

	move	t0, v0
	sw	t0,16($fp)
	li	v0, SYS_write 
	li	a0, 1         
	la	a1, msg   
	lw	t0,16($fp)
	move	a2, t0     
	syscall
	
	lw	$fp,28(sp)
	lw	ra,32(sp)
	lw	gp,24(sp)
	addu	sp,sp,40
	jr	ra
	.end	printHelp

	.rdata
msg:
	.asciiz	"Usage:\n	tp0 -h\n	tp0 -V\n	tp0 [file...]\nOptions:\n	-V, --version Print version and quit.\n	-h, --help Print this information and quit.\nExamples:\n	tp0 foo.txt bar.txt\n	tp0 gz.txt\n"
\end{lstlisting}

\subsection{printVersion.S}

\lstset{language=[x86masm]Assembler, breaklines=true, basicstyle=\normalsize}
\begin{lstlisting}
#include <mips/regdef.h>
#include <sys/syscall.h>

	.text
	.align	2
	.globl	printVersion
	.ent	printVersion
printVersion:
	.frame	$fp,40,ra	#4SRA + 2LTA + 4ABA
	.set	noreorder
	.cpload	t9
	.set	reorder
	subu 	sp,sp,40
	.cprestore 24
	sw	$fp,28(sp) 
	sw	ra,32(sp)
	move 	$fp, sp
	la 	a0, msg
	la	t9, strLength
	jal	ra, t9

	move	t0, v0
	sw	t0,16($fp)
	li	v0, SYS_write 
	li	a0, 1         
	la	a1, msg   
	lw	t0,16($fp)
	move	a2, t0     
	syscall
	
	lw	$fp,28(sp)
	lw	ra,32(sp)
	lw	gp,24(sp)
	addu	sp,sp,40
	jr	ra
	.end	printVersion

	.rdata
msg:
	.asciiz	"tp1 1.1\nCopyright © 2015 FIUBA.\nEsto es software libre: usted es libre de cambiarlo y redistribuirlo.\nNo hay NINGUNA GARANTÍA, hasta donde permite la ley.\n\nEscrito por Julian Scialabba, Kaoru Heanna y Paula Saffioti.\n" 
\end{lstlisting}

\subsection{isEndOfLine.S}

\lstset{language=[x86masm]Assembler, breaklines=true, basicstyle=\normalsize}
\begin{lstlisting}
#include <mips/regdef.h>
#include <sys/syscall.h>

	.text
	.align	2

	.globl	isEndOfLine
	.ent	isEndOfLine
isEndOfLine:
#creo frame
	.frame	$fp, 8, ra   # 2(SRA) + 0(FPA) + 0(LTA) + 0(ABA) (leaf function)
	.set	noreorder
	.cpload	t9
	.set	reorder
	subu	sp, sp, 8
# guardo registros
	.cprestore 0
	sw	$fp, 4(sp)
	move	$fp, sp
	sw 	a0, 8($fp) 		#guardo el primer argumento en el ABA de mi caller

	# Use v0 for the result.
	li	v0, 0
	# el codigo ascii de new line es 10
	li	t0, 10
	beq 	a0, t0, endOfLineTrue
	li	v0, 0
	j	isEndOfLine_return
endOfLineTrue:
	li	v0, 1
	j	isEndOfLine_return

isEndOfLine_return:
	# Destruimos el frame.
	# Restauro callee-saved regs
	lw	$fp, 4(sp)
	lw	gp, 0(sp)
	addu		sp, sp, 8

	# Retorno.
	#
	jr	ra
	.end	isEndOfLine
\end{lstlisting}

\subsection{strLength.S}

\lstset{language=[x86masm]Assembler, breaklines=true, basicstyle=\normalsize}
\begin{lstlisting}
#include <mips/regdef.h>
#include <sys/syscall.h>

	.text
	.align	2
	.globl	strLength
	.ent	strLength

strLength:
	.frame	$fp, 16, ra
	.set	noreorder
	.cpload	t9
	.set	reorder
        
        #creo stack frame
	subu	sp, sp, 16

        # Guardo $gp y $sp en SRA
	.cprestore 8
        sw	$fp, 12(sp)
        
        #De ahora en adelante utilizo $fp como sp
	move	$fp, sp

        # Guardo argumento en ABA de caller
        sw a0, 16($fp)

        # Creo variable incremental y la guardo en primer registro de LTA
        move t0, zero
        sw t0, 0($fp)

_for:
        # Condicion de corte del for
        addu t1, a0, t0  # Apunto en t1 al siguiente byte del string recibido por param
        lb t1, 0(t1) # Cargo el character apuntado en t1
        beq t1, zero, _end_for
        addu t0, t0, 1
        j _for

_end_for:
        
        # Preparo todo para retorno
        move v0, t0
        lw gp, 8($fp)
        lw $fp, 12(sp)
        addu sp, sp, 16
        jr ra

        .end strLength
\end{lstlisting}

\subsection{concatBuffer.S}

\lstset{language=[x86masm]Assembler, breaklines=true, basicstyle=\normalsize}
\begin{lstlisting}
#include <mips/regdef.h>
#include <sys/syscall.h>

# en LTA guardo: line, lineLength, buffer, bufferLength, ¿concatAddress?

	.text
	.align		2

	.globl		concatBuffer
	.ent		concatBuffer
concatBuffer:
#creo frame
	.frame		$fp, 56, ra   # 4(SRA) + 0(FPA) + 6(LTA) + 4(ABA) (leaf function)
	.set		noreorder
	.cpload		t9
	.set		reorder
	subu		sp, sp, 56

# guardo callee-saved registers
	.cprestore 	40 		# guardo gp
	sw		$fp, 44(sp)
	move		$fp, sp
	sw 		ra, 48($fp)
#guardo a0 y a1 en el ABA de mi caller
	sw 		a0, 56($fp)
	sw 		a1, 60($fp)

lineLength:
	move 		t0, a0			# t0 = line. pos 16
	sw 		t0, 16($fp)
	move 		t1, zero		# t1 = lineLength. pos 20
	sw 		t1, 20($fp)
	beq 		t0, zero, bufferLength 	#si line esta vacio, queda lineLength en 0
	move 		a0, t0
	la		t9, strLength
	jal		ra, t9
	move 		t1, v0			# guardo en t1 lo que devolvio strLength
	sw 		t1, 20($fp)

bufferLength:
	move 		t2, a1			# t2 = buffer. pos 24
	sw 		t2, 24($fp)
	move 		t3, zero		# t3 = bufferLength. pos 28
	sw 		t3, 28($fp)
	beq 		t2, zero, concat # si buffer esta vacio, queda bufferLength en 0
	move 		a0, t2
	la		t9, strLength
	jal		ra, t9
	move 		t3, v0			# guardo en t3 lo que devolvio strLength
	sw 		t3, 28($fp)
	bne 		t3, zero, concat
	lw 		v0, 16($fp)		# si bufferLength es 0, devuelvo line
	j		return

concat:
	lw		a0, 16($fp)		# string 1
	lw 		a1, 20($fp)		# len1
	lw 		a2, 24($fp)		# string 2
	lw 		a3, 28($fp)		# len2
	la		t9, mergeStrings
	jal		ra, t9
	j		return			# ya tengo guardado en v0 el string a retornar

return:
# Restauro callee-saved regs
	lw 		gp, 40(sp)
	lw 		$fp, 44(sp)
	lw 		ra, 48(sp)

# Destruimos el frame.
	addu		sp, sp, 56

# Retorno.
	jr		ra
	.end		concatBuffer
\end{lstlisting}

\subsection{printLines.S}

\lstset{language=[x86masm]Assembler, breaklines=true, basicstyle=\normalsize}
\begin{lstlisting}
#include <mips/regdef.h>
#include <sys/syscall.h>

	.text
	.align	2
	.globl	printLines
	.ent	printLines

printLines:
	.frame	$fp,40,ra	
	.set	noreorder
	.cpload	t9
	.set	reorder
	subu 	sp,sp,40
	.cprestore 24
	sw	$fp,28(sp) 
	sw	ra,32(sp)
	move 	$fp, sp

        sw      a0, 40($fp)
        sw      a1, 44($fp)
        move    t0, a0
        move    t1, a1
        subu    t0, t0, 1
        sw      t0, 16($fp)
        sll     t3, t0, 2   
        addu    t1, t1, t3  #Apunto a ultimo char*
        sw      t1, 20($fp)
_loop:  
        blt     t0, zero, _end_loop
        lw      t1, 20($fp)
        lw      t2, 0(t1)

        # Llamada strLength
	move 	a0, t2
	la	t9, strLength 
	jal	ra, t9
	move	t3, v0 #longitud

        # Llamo a syscall write
	li	v0, SYS_write 
	li	a0, 1   
	move 	a1, t2 #Contenido
	move	a2, t3    
	syscall

        # Libero memoria de char*
        lw      t1, 20($fp)
        lw 	a0, 0(t1)		
	la      t9, myfree
	jal	ra, t9

        # Actualizo estado variables
        lw      t0, 16($fp)
        lw      t1, 20($fp)
        subu    t0, t0, 1
        subu    t1, t1, 4
        sw      t0, 16($fp)
        sw      t1, 20($fp)

        j   _loop

_end_loop:
	lw	$fp,28(sp)
	lw	ra,32(sp)
	lw	gp,24(sp)
	addu	sp,sp,40
	jr	ra

	.end	printLines
\end{lstlisting}

\subsection{resizeArrayLines.S}

\lstset{language=[x86masm]Assembler, breaklines=true, basicstyle=\normalsize}
\begin{lstlisting}
#include <mips/regdef.h>
#include <sys/syscall.h>

	.text
	.align	2
	.globl	resizeArrayLines
	.ent	resizeArrayLines

resizeArrayLines:
	.frame	$fp,40,ra	
	.set	noreorder
	.cpload	t9
	.set	reorder
	subu 	sp,sp,40
	.cprestore 24
	sw	$fp,28(sp) 
	sw	ra,32(sp)
	move 	$fp, sp

        sw      a0, 40($fp)
        sw      a1, 44($fp)

        move    t0, zero
        sw      t0, 16($fp)

        addu	t1, a1, 1
	sll     t1, t1, 2
	move    a0, t1
	la	t9, mymalloc
	jal	ra, t9
	sw	v0, 20($fp)	

        lw      t0, 16($fp)
        lw      t1, 44($fp)
_loop:  
        bge     t0, t1, _end_loop

        lw      t2, 40($fp)
        lw      t3, 20($fp)

        sll     t4, t0, 2
        addu    t2, t2, t4
        addu    t3, t3, t4 
        lw      t4, 0(t2)
        sw      t4, 0(t3)

        addu    t0, t0, 1
        j   _loop

_end_loop:

        lw      a0, 40($fp)
	la      t9, myfree
	jal	ra, t9

        lw      v0, 20($fp)
	lw	$fp,28(sp)
	lw	ra,32(sp)
	lw	gp,24(sp)
	addu	sp,sp,40
	jr	ra

	.end	resizeArrayLines
\end{lstlisting}

\subsection{mergeStrings.S}

\lstset{language=[x86masm]Assembler, breaklines=true, basicstyle=\normalsize}
\begin{lstlisting}
#include <mips/regdef.h>
#include <sys/syscall.h>

# a0: char* string1 -> t0
# a1: size_t len1 -> t1
# a2: char* string2 -> t2
# a3: size_t len2 -> t3

	.text
	.align		2

	.globl		mergeStrings
	.ent		mergeStrings
mergeStrings:
#creo frame
	.frame		$fp, 56, ra   # 4(SRA) + 0(FPA) + 6(LTA) + 4(ABA) (leaf function)
	.set		noreorder
	.cpload		t9
	.set		reorder
	subu		sp, sp, 56

# guardo callee-saved registers
	.cprestore 	40 		# guardo gp
	sw		$fp, 44(sp)
	move		$fp, sp
	sw 		ra, 48($fp)
#guardo a0,a1,a2,a3 en el ABA de mi caller
	sw 		a0, 56($fp)
	sw 		a1, 60($fp)
	sw 		a2, 64($fp)
	sw 		a3, 68($fp)

	move 		t0, a0
	sw		t0, 16($fp)		# string1
	move 		t1, a1
	sw		t1, 20($fp)		# len1
	move 		t2, a2
	sw		t2, 24($fp)		# string2
	move 		t3, a3
	sw		t3, 28($fp)		# len2

allocMemory:
	addu		t4, t1, t3		# t4 = len1 + len2. pos 32
	sw		t4, 32($fp)
	addiu		t4, t4, 1		# t4 <- t4 +1
	move 		a0, t4			# cantidad de bytes a allocar
	la		t9, mymalloc
	jal		ra, t9
	sw		v0, 36($fp)		# dir(concat) = address allocated memory. pos 36

copyString1:
	move		t6, zero		# t6 = i
	lw		t0, 16($fp)		# t0 = &string1
	lw		t1, 20($fp)		# t1 = len1
	lw		t5, 36($fp)		# t5 = &concat
_forString1:
	bge 		t6, t1, copyString2		# si (i >= len1) voy a copyString2

	addu		t7, t0, t6		# t7 = &string1[i]
	lb 		t8, 0(t7) 		# t8 = string1[i]

	addu		t7, t5, t6		# t7 = &concat[i]
	sb		t8, 0(t7)		# concat[i] <- string1[i]

	addiu		t6, t6, 1		# i++
	j		_forString1

copyString2:
	move		t6, zero		# t6 = i
	lw		t2, 24($fp)		# t2 = &string2
	lw		t3, 28($fp)		# t3 = len2
	lw		t5, 36($fp)		# t5 = &concat
	lw		t1, 20($fp)		# t1 = len1
_forString2:
	bge 		t6, t3, endMerge		# si (i >= len2) voy a endMerge
	addu		t7, t2, t6		# t7 = &string2[i]
	lb		t8, 0(t7)		# t8 = string2[i]

	addu 		t7, t6, t1		# t7 = j <- i + len1
	addu		t7, t5, t7		# t7 = &concat[j]
	sb 		t8, 0(t7)		# concat[j] = string2[i]

	addiu		t6, t6, 1		#i++
	j		_forString2

endMerge:
	lw		t5, 36($fp)		# t5 = dir(concat)
	lw		t4, 32($fp)		# t4 = len1 + len2
	addu		t7, t5, t4		# t7 = &concat[len1+len2]
	sb		zero, 0(t7)		# concat[len1+len2] = 0
	lw		v0, 36($fp)		# devuelvo dir(concat)

return:
	# Restauro callee-saved regs
	lw 		gp, 40(sp)
	lw 		$fp, 44(sp)
	lw 		ra, 48(sp)

# Destruimos el frame.
	addu		sp, sp, 56

# Retorno.
	jr		ra
	.end		mergeStrings
\end{lstlisting}

\subsection{storeNewLine.S}

\lstset{language=[x86masm]Assembler, breaklines=true, basicstyle=\normalsize}
\begin{lstlisting}
#include <mips/regdef.h>
#include <sys/syscall.h>

	.text
	.align	2

	.globl	storeNewLine
	.ent	storeNewLine
storeNewLine:
	.frame	$fp,48,ra		

	.set	noreorder		
	.cpload	t9			
	.set	reorder			

	subu	sp,sp,48		

	.cprestore 32			

	sw	ra,36(sp)		
	sw	$fp,40(sp)		
	move	$fp,sp				

	sw	a0,20($fp)		
	sw	a1,24($fp)		
	sw	a2,16($fp)		
	sw	a3,28($fp)		
	
	lw	a0,20($fp)			
	la	t9,isEndOfLine		
	jal	ra,t9	
	
	beq	v0,zero,end		 

	lw	a0,16($fp)			

	lw	t0,24($fp)		
	lw	a1,0(t0)		
	la	t9,resizeArrayLines		
	jal	ra,t9
	sw	v0,16($fp)			

	lw	t0,24($fp)			
	lw	t0,0(t0)		
	lw	t1,16($fp)		

        sll     t3, t0, 2  	 	
        addu    t1, t1, t3  		

	lw	t0,28($fp)		
	lw	t2,0(t0)		
	sw	t2,0(t1)	
	li 	a0,1
	la	t9,mymalloc		
	jal	ra,t9	

	lw	t0,28($fp)		
	sw	v0,0(t0)		

	lw	t0,24($fp)			
	lw	t1,0(t0)		
	addu 	t1,t1,1			
	sw	t1,0(t0)	  		

	
end:
	lw	v0,16($fp)		
	lw	gp,32(sp)		
	lw	ra,36(sp)		
	lw	$fp,40(sp)		
	addu	sp,sp,48		
        jr 	ra

	.end	storeNewLine
	.rdata

\end{lstlisting}

\subsection{readFile.S}

\lstset{language=[x86masm]Assembler, breaklines=true, basicstyle=\normalsize}
\begin{lstlisting}
#include <mips/regdef.h>
#include <sys/syscall.h>

# a0 = buffer
# a1 = bufIncrSize
# a2 = filePointer

	.text
	.align	2
	.globl	readFile
	.ent	readFile

readFile:
	.frame	$fp, 32, ra
	.set	noreorder
	.cpload	t9
	.set	reorder
	subu	sp, sp, 32
	.cprestore 16
   	sw 	$fp, 20(sp)
    	sw 	ra, 24(sp)
    	move 	$fp,sp
   	sw	a0, 32($fp)
   	sw 	a1, 36($fp)
    	sw 	a2, 40($fp)
	
	jal fgets  
	
	lw 	gp, 16(sp)
	lw 	$fp, 20(sp)
	lw 	ra, 24(sp)
	addu 	sp,sp,32
	jr		ra
	.end readFile
\end{lstlisting}


\subsection{tacFile.S}

\lstset{language=[x86masm]Assembler, breaklines=true, basicstyle=\normalsize}
\begin{lstlisting}
#include <mips/regdef.h>
#include <sys/syscall.h>

	.text
	.align	2

	.globl	tacFile
	.ent	tacFile
tacFile:
	.frame	$fp,64,ra		

	.set	noreorder		
	.cpload	t9			
	.set	reorder			

	subu	sp,sp,64		

	.cprestore 48			

	sw	ra,56(sp)		
	sw	$fp,52(sp)		
	move	$fp,sp				
	
	sw	a0,64($fp)		

	li	t0,0			
	sw 	t0,16($fp)		
	li 	a0,1
	la	t9,mymalloc		
	jal	ra,t9		
	sw	v0,20($fp)		

	li 	a0,10
	la	t9,mymalloc	
	jal	ra,t9		
	sw	v0,36($fp)		


	li	a0,1			
	la	t9,mymalloc		
	jal	ra,t9		
	sw	v0,24($fp)			

cond_while: 

	lw	a0,36($fp)
	li	a1,10
	lw	a2,64($fp)
	la	t9,readFile
	jal	ra,t9

	bne	v0,zero,bloque_while	
	b	end

bloque_while: 
	sw	v0,36($fp)             

	lw	a0,24($fp)		
	lw	a1,36($fp)		
	la	t9,concatBuffer		
	jal	ra,t9
	sw	v0,28($fp)		

	lw	a0,24($fp)		
	la	t9,myfree		
	jal	ra,t9	

	lw	t0,28($fp)		
	sw	t0,24($fp)	

 	lw      a0, 36($fp)		
	la	t9,strLength		
	jal	ra,t9
	move 	t0,v0	
	subu	t0,t0,1
	lw	t1,36($fp)
	addu	t2,t1,t0	
	lb	t3,0(t2)
	sw	t3,32($fp)		

	lw	a0,32($fp)		
	addu	a1,$fp,16		
	lw	a2,20($fp)		
	addu	a3,$fp,24		
	la	t9,storeNewLine		
	jal	ra,t9	
	sw	v0,20($fp)
	
	b	cond_while
end: 

	lw	a0,16($fp)		
	lw	a1,20($fp)		
	la	t9,printLines		
	jal	ra,t9

	lw	a0,24($fp)		
	la	t9,myfree		
	jal	ra,t9

	lw	a0,20($fp)		
	la	t9,myfree		
	jal	ra,t9

  	lw	a0,36($fp)		
	la	t9,myfree		
	jal	ra,t9

	li 	v0,1

	lw	gp,48(sp)		
	lw	ra,56(sp)		
	lw	$fp,52(sp)		
	addu	sp,sp,64		
        jr 	ra

	.end	tacFile

	.rdata
	.align 2
\end{lstlisting}

\section{Enunciado}

\maketitle
\end{document}
